% ==========================================
% BAB VI EVALUASI
% ==========================================
\cleardoublepage
\chapter{EVALUASI}
\label{chap:evaluasi}

\section{Metode Evaluasi}

Evaluasi Aplikasi Pembelajaran ASD dilakukan untuk memastikan bahwa sistem yang dikembangkan telah memenuhi seluruh kebutuhan fungsional yang ditetapkan, menghasilkan konten pembelajaran yang berkualitas, serta dapat digunakan dengan mudah oleh mahasiswa sebagai pengguna. Evaluasi terdiri atas dua bagian utama, yaitu pengujian \textit{black box} untuk memverifikasi fungsionalitas sistem dan pengujian kegunaan (\textit{usability testing}) untuk menilai pengalaman pengguna secara keseluruhan.

\subsection{Pengujian \textit{Black Box}}

\textit{Black Box Testing} merupakan metode pengujian perangkat lunak yang mengevaluasi fungsi dan perilaku sistem dari sisi eksternal tanpa melihat struktur internal atau kode program, sehingga pengujian dilakukan dengan memberikan berbagai \textit{input} dan memeriksa apakah \textit{output} telah sesuai dengan spesifikasi yang ditentukan \cite{Pressman2015SWE}. Pada pengujian ini digunakan teknik \textit{Equivalence Partitioning}, yaitu metode yang membagi domain \textit{input} ke dalam beberapa kelas ekivalen sehingga hanya satu nilai dari setiap kelas yang perlu diuji karena dianggap mewakili seluruh anggota kelas tersebut \cite{Myers2011ArtOfSoftwareTesting}.

Pengujian \textit{black box} dilakukan terhadap seluruh fitur fungsional Aplikasi Pembelajaran ASD berdasarkan kebutuhan fungsional yang telah ditetapkan. Setiap fitur diuji dengan mendefinisikan skenario uji, data masukan, keluaran yang diharapkan, dan status kelulusan. Rangkuman cakupan pengujian \textit{black box} dapat dilihat pada Tabel \ref{tab:blackbox} berikut.

\begin{longtable}{|L{0.9cm}|L{2.3cm}|L{4.4cm}|L{4.4cm}|}
\caption{Cakupan Pengujian \textit{Black Box}}
\label{tab:blackbox}\\
\hline
\textbf{ID} & \textbf{Fitur} & \textbf{Skenario Uji} & \textbf{Indikator Keberhasilan} \\
\hline
\endfirsthead

\caption[]{Cakupan Pengujian \textit{Black Box} (lanjutan)}\\
\hline
\textbf{ID} & \textbf{Fitur} & \textbf{Skenario Uji} & \textbf{Indikator Keberhasilan} \\
\hline
\endhead

\hline
\multicolumn{4}{r}{\textit{Bersambung ke halaman berikutnya}}\\
\endfoot

\hline
\endlastfoot

FT-01 & Pilih Topik ASD &
Pengguna memilih salah satu dari sepuluh topik yang tersedia &
Sistem menampilkan halaman topik dengan tab sesuai status \textit{guided flow} pengguna \\
\hline

FT-02 & Tampilkan Materi &
Pengguna membuka tab materi pada topik yang dipilih &
Konten materi statis berhasil dirender dan ditampilkan sesuai topik, progres diperbarui setelah selesai dibaca \\
\hline

FT-03 & Tampilkan Contoh &
Pengguna membuka tab contoh setelah menyelesaikan tab materi &
Konten contoh berhasil dihasilkan dan ditampilkan, tab terkunci jika materi belum selesai \\
\hline

FT-04 & Latihan Soal &
Pengguna membuka tab latihan dan menerima soal yang di-\textit{generate} &
Soal latihan berhasil di-\textit{generate} oleh LLM dan ditampilkan kepada pengguna \\
\hline

FT-05 & Pembahasan Latihan &
Pengguna mengumpulkan jawaban latihan &
Sistem mengembalikan skor, \textit{feedback} per soal, dan pembahasan yang relevan \\
\hline

FT-06 & Ringkasan Materi &
Pengguna membuka tab ringkasan setelah menyelesaikan latihan &
Konten ringkasan statis berhasil dirender dan ditampilkan, topik ditandai selesai \\
\hline

FT-07 & Progres Belajar &
Pengguna membuka halaman progres belajar &
Sistem menampilkan status penyelesaian seluruh topik beserta indikator kelemahan konsep \\
\hline

FT-08 & \textit{Generate} Konten LLM &
Sistem mengirim \textit{prompt} ke OpenAI API untuk menghasilkan materi, contoh, dan soal &
OpenAI API mengembalikan konten dalam format JSON yang valid sesuai topik \\
\hline

FT-09 & \textit{Guided Learning Flow} &
Pengguna mencoba mengakses tab yang belum terbuka &
Sistem mengunci tab yang belum tersedia dan menampilkan notifikasi untuk menyelesaikan tahap sebelumnya \\
\hline

FT-10 & \textit{Feedback} Jawaban &
Pengguna mengumpulkan jawaban latihan &
Sistem menampilkan \textit{feedback} spesifik per soal beserta skor yang akurat \\
\hline

FT-11 & \textit{Adaptive Learning Engine} &
Pengguna menyelesaikan latihan dengan hasil tertentu &
Sistem menyimpan \textit{weak concepts} dan menggunakannya pada \textit{generate} soal berikutnya \\
\hline

FT-12 & Adaptasi Konten Pembelajaran &
Pengguna dengan riwayat \textit{weak concepts} membuka materi/contoh/latihan &
Konten yang dihasilkan lebih menekankan konsep yang teridentifikasi sebagai kelemahan pengguna \\
\hline

FT-13 & Evaluasi Soal Sendiri &
Pengguna memasukkan soal dan jawaban dari sumber eksternal pada panel Soal Sendiri. Skenario uji meliputi soal berisi teks, soal diunggah dari file \texttt{.txt}, dan kolom jawaban dikosongkan &
Sistem mengembalikan hasil evaluasi berisi skor, nilai, komentar, bagian benar/salah, dan rekomendasi belajar. Soal tanpa jawaban mendapat nilai "Belum Dijawab" \\
\hline

FT-14 & Chat Asisten Materi &
Pengguna mengirim pertanyaan melalui \textit{widget chat} pada halaman topik aktif, skenario uji meliputi pertanyaan relevan dengan topik dan pertanyaan lanjutan yang mengacu pada jawaban sebelumnya &
Sistem mengembalikan balasan kontekstual yang relevan dengan materi topik, pertanyaan lanjutan dijawab dengan mempertimbangkan riwayat percakapan \\
\hline

\end{longtable}


\subsection{Pengujian Kegunaan (\textit{Usability Testing})}

Selain pengujian fungsional, dilakukan pengujian kegunaan untuk menilai sejauh mana aplikasi dapat digunakan oleh mahasiswa dengan mudah, efisien, dan memuaskan. Pengujian kegunaan dilakukan melalui dua pendekatan, yaitu \textit{task-based usability test} dan pengisian kuesioner \textit{System Usability Scale} (SUS).

\subsubsection{\textit{Task-Based Usability Test}}

\textit{Task-based usability test} dilakukan dengan meminta responden menyelesaikan serangkaian tugas yang mencerminkan alur penggunaan aplikasi secara nyata. Responden diminta untuk menyelesaikan tugas-tugas berikut secara berurutan tanpa bantuan dari peneliti:

\begin{enumerate}
    \item Memilih salah satu topik ASD dari halaman \textit{dashboard}.
    \item Membaca konten materi dan menandai tab materi sebagai selesai.
    \item Membuka tab contoh dan membaca contoh implementasi yang ditampilkan.
    \item Mengerjakan latihan soal dan mengumpulkan seluruh jawaban.
    \item Membaca pembahasan dan skor yang diberikan sistem.
    \item Membuka tab ringkasan dan menyelesaikan alur belajar satu topik.
    \item Mengakses halaman progres belajar dan melihat status penyelesaian topik.
\end{enumerate}

Selama pengujian, peneliti mengamati dan mencatat waktu penyelesaian setiap tugas, kesalahan yang dilakukan, serta hambatan yang ditemui. Indikator keberhasilan \textit{task-based usability test} adalah seluruh tugas dapat diselesaikan oleh responden dengan waktu yang wajar tanpa adanya kebingungan yang signifikan dalam mengikuti alur \textit{guided learning}.

\subsubsection{\textit{System Usability Scale} (SUS)}

\textit{System Usability Scale} (SUS) merupakan instrumen pengukuran kegunaan yang dikembangkan oleh John Brooke (1996) dan telah banyak digunakan dalam evaluasi sistem interaktif. SUS terdiri atas 10 pernyataan dengan skala Likert 1–5, di mana skor akhir dinormalisasi ke rentang 0–100. Interpretasi skor SUS mengacu pada ketentuan yang dapat dilihat pada Tabel \ref{tab:sus}.

\begin{table}[H]
\centering
\caption{Interpretasi Skor SUS}
\label{tab:sus}
\begin{tabular}{|L{3cm}|L{3cm}|L{6cm}|}
\hline
\textbf{Rentang Skor} & \textbf{Kategori} & \textbf{Keterangan} \\
\hline
$> 80{,}3$ & \textit{Excellent} & Aplikasi sangat mudah digunakan \\
\hline
$68 - 80{,}3$ & \textit{Good} & Aplikasi mudah digunakan \\
\hline
$51 - 68$ & \textit{OK} & Aplikasi cukup dapat digunakan \\
\hline
$< 51$ & \textit{Poor} & Aplikasi sulit digunakan dan perlu perbaikan signifikan \\
\hline
\end{tabular}
\end{table}

Kuesioner SUS disebarkan kepada responden melalui Google Form setelah responden menyelesaikan \textit{task-based usability test}. Target responden adalah mahasiswa yang sedang atau telah menempuh mata kuliah Algoritma dan Struktur Data. Indikator keberhasilan pengujian kegunaan adalah perolehan skor SUS minimal 68 yang menunjukkan bahwa aplikasi berada pada kategori \textit{Good} dan layak digunakan sebagai media pembelajaran.

\subsubsection{Responden}

Responden pengujian kegunaan merupakan mahasiswa yang sedang atau telah menempuh mata kuliah Algoritma dan Struktur Data. Kriteria responden ditetapkan sebagai berikut:

\begin{enumerate}
    \item Mahasiswa aktif yang sedang atau telah menempuh mata kuliah Algoritma dan Struktur Data.
    \item Memiliki pengalaman menggunakan perangkat komputer atau \textit{smartphone} untuk keperluan akademik.
    \item Bersedia mengikuti seluruh tahapan pengujian secara mandiri.
\end{enumerate}

Pengujian dilakukan terhadap minimal 5 responden sesuai dengan panduan Nielsen (1993) yang menyatakan bahwa 5 pengguna sudah cukup untuk mengidentifikasi sebagian besar permasalahan kegunaan dalam sebuah antarmuka.

\subsection{Evaluasi Kualitas Konten LLM}

Selain pengujian fungsional dan kegunaan, dilakukan evaluasi terhadap kualitas konten yang dihasilkan oleh LLM untuk menjawab rumusan masalah kedua, yaitu sejauh mana LLM mampu menghasilkan konten yang relevan, akurat, dan konsisten dengan kurikulum ASD. Evaluasi ini difokuskan pada dua jenis konten yang dihasilkan secara dinamis oleh sistem: soal latihan dan evaluasi jawaban pengguna.

\subsubsection{Metodologi}

Evaluasi kualitas konten dilakukan dengan metode \textit{expert-based content validation}, yaitu membandingkan keluaran LLM terhadap jawaban referensi yang telah divalidasi terlebih dahulu berdasarkan materi kurikulum ASD. Proses evaluasi dilakukan dalam tiga tahap:

\begin{enumerate}
    \item \textbf{Penyusunan kisi-kisi soal referensi.} Disiapkan soal-soal uji beserta jawaban referensi yang benar berdasarkan sumber kurikulum (buku teks ASD dan modul kuliah). Soal mencakup tiga topik representatif yang mewakili tingkat kesulitan berbeda: \textit{Linked List} (dasar), \textit{Tree} (menengah), dan \textit{Graph} (lanjutan), masing-masing terdiri atas 5 soal sehingga total terdapat 15 soal referensi.

    \item \textbf{Pengujian pada sistem.} Setiap soal referensi dimasukkan ke dalam fitur Soal Sendiri (\textit{evaluasi soal mandiri}) pada aplikasi bersama jawaban referensinya. Sistem menghasilkan hasil evaluasi berupa skor, komentar, dan umpan balik melalui LLM.

    \item \textbf{Penilaian kesesuaian.} Keluaran LLM dinilai secara manual menggunakan rubrik tiga kriteria sebagaimana tercantum pada Tabel \ref{tab:rubrik-llm}. Setiap kriteria diberi bobot dan skor 1–3 (tidak sesuai / sebagian sesuai / sesuai), kemudian dikonversi ke persentase kesesuaian.
\end{enumerate}

\begin{table}[H]
\centering
\caption{Rubrik Evaluasi Kualitas Konten LLM}
\label{tab:rubrik-llm}
\begin{tabular}{|L{0.5cm}|L{3.5cm}|L{1.2cm}|L{7cm}|}
\hline
\textbf{No} & \textbf{Kriteria} & \textbf{Bobot} & \textbf{Deskripsi} \\
\hline
1 & Kebenaran Konseptual & 40\% & Penjelasan atau penilaian LLM secara konseptual benar sesuai prinsip ASD yang berlaku \\
\hline
2 & Kesesuaian Kurikulum & 30\% & Konten merujuk pada terminologi, notasi, dan cakupan topik yang sesuai dengan kurikulum ASD yang diajarkan \\
\hline
3 & Kejelasan dan Tidak Menyesatkan & 30\% & Penjelasan disampaikan dengan jelas, tidak ambigu, dan tidak mengandung informasi yang dapat menimbulkan miskonsepsi \\
\hline
\end{tabular}
\end{table}

\subsubsection{Indikator Keberhasilan}

Indikator keberhasilan evaluasi kualitas konten LLM adalah perolehan rata-rata persentase kesesuaian konten minimal 70\% pada seluruh kriteria. Nilai ambang ini dipilih berdasarkan pertimbangan bahwa konten dengan tingkat kesesuaian di atas 70\% dianggap layak digunakan sebagai referensi pembelajaran pendukung tanpa menimbulkan miskonsepsi yang signifikan.


\section{Hasil Evaluasi}

Bagian ini menyajikan hasil dari tiga tahap evaluasi yang telah dilakukan terhadap aplikasi Teman Belajar ASD, yaitu pengujian \textit{black box} untuk memverifikasi pemenuhan kebutuhan fungsional, pengujian kegunaan untuk menilai kemudahan penggunaan aplikasi oleh mahasiswa, serta evaluasi kualitas konten LLM untuk mengukur akurasi dan relevansi keluaran model terhadap kurikulum.

\subsection{Hasil Pengujian \textit{Black Box}}

Pengujian \textit{black box} dilakukan terhadap 14 fitur fungsional aplikasi dengan total 18 skenario uji menggunakan teknik \textit{Equivalence Partitioning}. Setiap skenario diuji dengan data masukan yang representatif dari masing-masing kelas ekivalen. Hasil pengujian disajikan pada Tabel \ref{tab:hasil-blackbox} berikut.

\begin{longtable}{|L{0.8cm}|L{1.8cm}|L{2.2cm}|L{2.5cm}|L{2.5cm}|L{1.2cm}|}
\caption{Hasil Pengujian \textit{Black Box}}
\label{tab:hasil-blackbox}\\
\hline
\textbf{ID} & \textbf{Skenario Uji} & \textbf{Data Masukan} & \textbf{Keluaran yang Diharapkan} & \textbf{Keluaran Aktual} & \textbf{Status} \\
\hline
\endfirsthead

\caption[]{Hasil Pengujian \textit{Black Box} (lanjutan)}\\
\hline
\textbf{ID} & \textbf{Skenario Uji} & \textbf{Data Masukan} & \textbf{Keluaran yang Diharapkan} & \textbf{Keluaran Aktual} & \textbf{Status} \\
\hline
\endhead

\hline
\multicolumn{6}{r}{\textit{Bersambung ke halaman berikutnya}}\\
\endfoot

\hline
\endlastfoot

FT-01 & Memilih topik ASD &
Pengguna mengklik salah satu dari sepuluh kartu topik pada halaman topik &
Sistem menampilkan halaman pembelajaran topik dengan tab sesuai status \textit{guided flow} pengguna &
Halaman topik berhasil dimuat dengan tab yang sesuai status pengguna &
Lulus \\
\hline

FT-02 & Men\-ampilkan materi &
Pengguna membuka tab materi pada topik yang dipilih &
Konten materi statis berhasil dirender dan ditampilkan, progres tab materi diperbarui setelah selesai dibaca &
Materi berhasil ditampilkan dan progres diperbarui setelah pengguna menyelesaikan tab &
Lulus \\
\hline

FT-03a & Membuka tab contoh setelah materi selesai &
Pengguna mengklik tab contoh setelah menandai materi sebagai selesai &
Tab contoh terbuka dan konten contoh beserta kode berhasil ditampilkan &
Tab contoh berhasil terbuka dan konten ditampilkan &
Lulus \\
\hline

FT-03b & Membuka tab contoh sebelum materi selesai &
Pengguna mencoba mengklik tab contoh tanpa menyelesaikan materi &
Tab terkunci dan sistem menampilkan notifikasi untuk menyelesaikan materi terlebih dahulu &
Tab terkunci dan notifikasi ditampilkan sesuai yang diharapkan &
Lulus \\
\hline

FT-04 & \textit{Generate} soal latihan &
Pengguna membuka tab latihan setelah menyelesaikan tab contoh &
LLM menghasilkan 5 soal latihan yang relevan dengan topik yang sedang dipelajari &
5 soal latihan berhasil di-\textit{generate} dan ditampilkan kepada pengguna &
Lulus \\
\hline

FT-05 & Meng\-umpulkan jawaban latihan &
Pengguna menjawab seluruh soal dan menekan tombol kumpulkan &
Sistem mengembalikan skor keseluruhan, \textit{feedback} per soal, dan konsep yang perlu diperbaiki &
Skor, \textit{feedback} per soal, dan konsep lemah berhasil ditampilkan &
Lulus \\
\hline

FT-06 & Membuka tab ringkasan &
Pengguna membuka tab ringkasan setelah menyelesaikan latihan &
Konten ringkasan statis berhasil dirender dan ditampilkan, topik ditandai sebagai selesai 100\% &
Ringkasan berhasil ditampilkan dan status topik diperbarui menjadi selesai &
Lulus \\
\hline

FT-07 & Melihat progres belajar &
Pengguna mengakses halaman progres &
Sistem menampilkan status penyelesaian seluruh topik, \textit{streak} belajar, dan rata-rata progres &
Halaman progres berhasil dimuat dengan data lengkap sesuai riwayat pengguna &
Lulus \\
\hline

FT-08 & \textit{Generate} konten LLM &
Sistem mengirim \textit{prompt} berstruktur ke OpenAI API &
OpenAI API mengembalikan konten dalam format JSON yang valid dan sesuai topik &
Respons JSON valid diterima dan berhasil di-\textit{parse} untuk ditampilkan ke pengguna &
Lulus \\
\hline

FT-09 & Akses tab yang terkunci &
Pengguna mencoba mengakses tab yang belum tersedia dalam alur \textit{guided learning} &
Sistem mengunci tab dan menampilkan notifikasi untuk menyelesaikan tahap sebelumnya &
Tab terkunci dan notifikasi muncul sesuai yang diharapkan &
Lulus \\
\hline

FT-10 & \textit{Feedback} per soal &
Pengguna mengumpul\-kan jawaban latihan yang beragam (benar, salah, sebagian benar) &
Sistem menampilkan \textit{feedback} spesifik untuk setiap soal beserta penjelasan jawaban yang benar &
\textit{Feedback} spesifik per soal berhasil ditampilkan dengan penjelasan yang relevan &
Lulus \\
\hline

FT-11 & \textit{Adaptive learning engine} &
Pengguna menyelesaikan latihan dengan beberapa konsep yang belum dikuasai &
Sistem menyimpan daftar \textit{weak concepts} dan meneruskannya sebagai parameter pada \textit{generate} soal berikutnya &
\textit{Weak concepts} berhasil disimpan dan digunakan sebagai konteks pada \textit{generate} soal sesi berikutnya &
Lulus \\
\hline

FT-12 & Adaptasi konten ber\-dasarkan kelemahan &
Pengguna dengan riwayat \textit{weak concepts} membuka fitur \textit{generate} soal baru &
Soal yang dihasilkan lebih menekankan konsep yang sebelumnya teridentifikasi sebagai kelemahan pengguna &
Soal yang dihasilkan mencakup konsep-konsep lemah pengguna sesuai riwayat sebelumnya &
Lulus \\
\hline

FT-13a & Evaluasi soal sendiri dengan teks jawaban &
Pengguna memasukkan soal dan jawaban pada panel Soal Sendiri lalu menekan tombol evaluasi &
Sistem mengembalikan skor, nilai, komentar, bagian benar, bagian perlu diperbaiki, dan rekomendasi belajar &
Hasil evaluasi berhasil ditampilkan lengkap sesuai yang diharapkan &
Lulus \\
\hline

FT-13b & Evaluasi soal sendiri tanpa jawaban &
Pengguna memasukkan soal tanpa mengisi kolom jawaban &
Sistem mengembalikan nilai "Belum Dijawab" dengan skor 0 &
Sistem menampilkan nilai "Belum Dijawab" sesuai yang diharapkan &
Lulus \\
\hline

FT-14a & Chat asisten dengan pertanyaan relevan &
Pengguna mengirim pertanyaan seputar materi topik aktif melalui \textit{widget chat} &
Sistem mengembalikan balasan kontekstual yang relevan dengan materi topik tersebut &
Balasan berhasil ditampilkan dan relevan dengan konteks materi &
Lulus \\
\hline

FT-14b & Chat asisten dengan pertanyaan lanjutan &
Pengguna mengirim pertanyaan lanjutan yang mengacu pada jawaban sebelumnya dalam satu sesi &
Sistem menjawab dengan mempertimbangkan riwayat percakapan sehingga konteks tetap terjaga &
Balasan mempertimbangkan riwayat percakapan sesuai yang diharapkan &
Lulus \\
\hline

\end{longtable}

Berdasarkan hasil pengujian pada Tabel \ref{tab:hasil-blackbox}, seluruh 18 skenario uji dari 14 fitur fungsional dinyatakan \textbf{lulus}. Hasil ini menunjukkan bahwa aplikasi Teman Belajar ASD telah memenuhi seluruh kebutuhan fungsional yang ditetapkan pada Tabel \ref{tab:kebutuhan-fungsional}, mulai dari alur \textit{guided learning}, integrasi LLM, hingga mekanisme pembelajaran adaptif.

\subsection{Hasil Pengujian Kegunaan}

Pengujian kegunaan dilakukan terhadap 5 responden mahasiswa yang sedang atau telah menempuh mata kuliah Algoritma dan Struktur Data, sesuai kriteria yang telah ditetapkan pada bagian sebelumnya. Profil responden disajikan pada Tabel \ref{tab:profil-responden}.

\begin{table}[H]
\centering
\caption{Profil Responden Pengujian Kegunaan}
\label{tab:profil-responden}
\begin{tabular}{|L{1cm}|L{3.6cm}|L{4cm}|L{1.6cm}|L{1.2cm}|}
\hline
\textbf{Kode} & \textbf{Nama} & \textbf{Asal Institusi} & \textbf{Angkatan} & \textbf{Umur} \\
\hline
R1 & Muhammad Zufar Dafy & Institut Teknologi Nasional & 2022 & 22 \\
\hline
R2 & Amira Ghina Nurfansepta & Universitas Brawijaya & 2021 & 22 \\
\hline
R3 & Muhammad Faiz Fahri & Universitas Padjadjaran & 2022 & 21 \\
\hline
R4 & Thea & Institut Teknologi Bandung & 2022 & 21 \\
\hline
R5 & Vieri & Institut Teknologi Bandung & 2021 & 24 \\
\hline
\end{tabular}
\end{table}

Seluruh responden merupakan mahasiswa yang sudah pernah menempuh mata kuliah Algoritma dan Struktur Data dan baru pertama kali menggunakan aplikasi Teman Belajar ASD.

\subsubsection{Hasil \textit{Task-Based Usability Test}}

Kelima responden diminta menyelesaikan tujuh tugas (T1-T7) pada topik \textit{Linked List} secara mandiri tanpa bantuan peneliti, sambil diamati waktu penyelesaian dan hambatan yang dialami. Rekapitulasi hasil observasi disajikan pada Tabel \ref{tab:hasil-usability-task}.

\begin{table}[H]
\centering
\caption{Rekapitulasi Hasil \textit{Task-Based Usability Test}}
\label{tab:hasil-usability-task}
\begin{tabular}{|L{3.6cm}|L{2.2cm}|L{2.6cm}|L{2.6cm}|}
\hline
\textbf{Tugas} & \textbf{Jumlah Selesai} & \textbf{Tingkat Keberhasilan} & \textbf{Rata-rata Waktu (detik)} \\
\hline
T1 - Pilih Topik & 5/5 & 100\% & 10,2 \\
\hline
T2 - Materi & 5/5 & 100\% & 100,6 \\
\hline
T3 - Contoh & 5/5 & 100\% & 53,0 \\
\hline
T4 - Latihan & 5/5 & 100\% & 206,2 \\
\hline
T5 - Pembahasan & 5/5 & 100\% & 35,6 \\
\hline
T6 - Ringkasan & 5/5 & 100\% & 30,2 \\
\hline
T7 - Progres & 5/5 & 100\% & 12,4 \\
\hline
\textbf{Rata-rata Keseluruhan} & \textbf{5/5} & \textbf{100\%} & \textbf{64,0} \\
\hline
\end{tabular}
\end{table}

Berdasarkan Tabel \ref{tab:hasil-usability-task}, seluruh responden berhasil menyelesaikan ketujuh tugas dengan tingkat keberhasilan 100\% tanpa ada yang menyerah di tengah pengujian. Tugas T4 (Latihan) memiliki rata-rata waktu penyelesaian paling lama (206,2 detik) karena mencakup waktu tunggu proses \textit{generate} soal oleh LLM serta waktu pengerjaan dan pembacaan evaluasi, sedangkan T1 (Pilih Topik) dan T7 (Progres) paling cepat diselesaikan (di bawah 15 detik) karena hanya melibatkan satu kali navigasi.

Meskipun seluruh tugas formal tercatat tanpa hambatan yang menghentikan penyelesaian, sesi \textit{think-aloud} dan wawancara pasca-tugas mengungkap beberapa catatan kegunaan yang bersifat minor:
\begin{itemize}
    \item Waktu tunggu \textit{generate} soal oleh LLM dirasakan cukup lama oleh beberapa responden (R1, R2, R5), meskipun dipahami sebagai konsekuensi dari pemrosesan AI.
    \item Label tab "Contoh" sempat menimbulkan sedikit kebingungan awal karena dikira berupa deskripsi materi, bukan soal contoh beserta pembahasannya (R3).
    \item Kotak jawaban pengguna dan kotak \textit{feedback} "jawaban tidak lengkap" pada fitur latihan menggunakan warna abu-abu yang serupa sehingga kurang mudah dibedakan (R4).
    \item Validasi/penilaian jawaban oleh AI pada satu kasus dirasa terlalu ketat oleh salah satu responden (R5), yang berpendapat evaluasi seharusnya bisa dilakukan manusia untuk kasus tertentu.
    \item Ukuran teks pada navigasi dirasa relatif kecil (R1).
    \item Sebagian responden mengharapkan tambahan visualisasi/ilustrasi grafis untuk materi struktur data, terutama \textit{Linked List} (R3, R4).
\end{itemize}

Seluruh catatan tersebut bersifat perbaikan pengalaman pengguna (\textit{enhancement}) dan tidak menghalangi responden menyelesaikan alur \textit{guided learning} secara keseluruhan.

\subsubsection{Hasil \textit{System Usability Scale} (SUS)}

Setelah menyelesaikan \textit{task-based usability test}, kelima responden mengisi kuesioner SUS. Skor tiap responden dihitung sesuai ketentuan standar SUS (butir ganjil: skor jawaban dikurangi 1; butir genap: 5 dikurangi skor jawaban; jumlah seluruh butir dikali 2,5). Hasil perhitungan skor SUS disajikan pada Tabel \ref{tab:hasil-sus}.

\begin{table}[H]
\centering
\caption{Hasil Skor \textit{System Usability Scale} (SUS)}
\label{tab:hasil-sus}
\begin{tabular}{|L{1cm}|L{3.5cm}|L{2.5cm}|L{3cm}|}
\hline
\textbf{Kode} & \textbf{Responden} & \textbf{Skor SUS} & \textbf{Kategori} \\
\hline
R1 & Responden 1 & 87,5 & \textit{Excellent} \\
\hline
R2 & Responden 2 & 97,5 & \textit{Excellent} \\
\hline
R3 & Responden 3 & 85,0 & \textit{Excellent} \\
\hline
R4 & Responden 4 & 95,0 & \textit{Excellent} \\
\hline
R5 & Responden 5 & 100,0 & \textit{Excellent} \\
\hline
\textbf{Rata-rata} & & \textbf{93,0} & \textbf{\textit{Excellent}} \\
\hline
\end{tabular}
\end{table}

Rata-rata skor SUS yang diperoleh adalah 93,0, yang berdasarkan Tabel \ref{tab:sus} berada pada kategori \textit{Excellent} ($> 80{,}3$). Seluruh responden secara individu juga memperoleh skor pada kategori \textit{Excellent}, dengan skor terendah 85,0 (R3) dan skor tertinggi 100,0 (R5). Hasil ini jauh melampaui indikator keberhasilan minimal, yaitu skor SUS $\geq 68$ (kategori \textit{Good}), sehingga dapat disimpulkan bahwa aplikasi Teman Belajar ASD sangat mudah digunakan oleh mahasiswa sebagai media pembelajaran mandiri.

\subsection{Hasil Evaluasi Kualitas Konten LLM}

Evaluasi kualitas konten LLM dilakukan terhadap 15 soal referensi yang terbagi ke dalam tiga topik, yaitu \textit{Linked List}, \textit{Tree}, dan \textit{Graph}, masing-masing 5 soal, menggunakan fitur Soal Sendiri pada aplikasi. Setiap keluaran evaluasi LLM dinilai menggunakan rubrik pada Tabel \ref{tab:rubrik-llm} dengan skor 1–3 untuk tiap kriteria. Hasil penilaian selengkapnya disajikan pada Tabel \ref{tab:hasil-llm}.

\begin{longtable}{|L{0.6cm}|L{1.6cm}|L{1.8cm}|L{1.3cm}|L{1.3cm}|L{1.3cm}|L{1.8cm}|}
\caption{Hasil Penilaian Kualitas Konten LLM}
\label{tab:hasil-llm}\\
\hline
\textbf{No} & \textbf{Kode Soal} & \textbf{Topik} & \textbf{K1} & \textbf{K2} & \textbf{K3} & \textbf{Skor (\%)} \\
\hline
\endfirsthead

\caption[]{Hasil Penilaian Kualitas Konten LLM (lanjutan)}\\
\hline
\textbf{No} & \textbf{Kode Soal} & \textbf{Topik} & \textbf{K1} & \textbf{K2} & \textbf{K3} & \textbf{Skor (\%)} \\
\hline
\endhead

\hline
\multicolumn{7}{r}{\textit{Bersambung ke halaman berikutnya}}\\
\endfoot

\hline
\endlastfoot

1 & LL-01 & Linked List & 3 & 3 & 3 & 100 \\
\hline
2 & LL-02 & Linked List & 3 & 3 & 3 & 100 \\
\hline
3 & LL-03 & Linked List & 3 & 3 & 3 & 100 \\
\hline
4 & LL-04 & Linked List & 3 & 3 & 2 & 90 \\
\hline
5 & LL-05 & Linked List & 3 & 3 & 3 & 100 \\
\hline
6 & TR-01 & Tree & 3 & 3 & 2 & 90 \\
\hline
7 & TR-02 & Tree & 3 & 3 & 3 & 100 \\
\hline
8 & TR-03 & Tree & 3 & 3 & 2 & 90 \\
\hline
9 & TR-04 & Tree & 3 & 3 & 2 & 90 \\
\hline
10 & TR-05 & Tree & 3 & 3 & 3 & 100 \\
\hline
11 & GR-01 & Graph & 3 & 3 & 3 & 100 \\
\hline
12 & GR-02 & Graph & 3 & 3 & 3 & 100 \\
\hline
13 & GR-03 & Graph & 3 & 3 & 3 & 100 \\
\hline
14 & GR-04 & Graph & 3 & 3 & 3 & 100 \\
\hline
15 & GR-05 & Graph & 3 & 3 & 3 & 100 \\
\hline
\textbf{Rata-rata} & & & \textbf{3,00} & \textbf{3,00} & \textbf{2,73} & \textbf{97,33} \\
\hline
\end{longtable}

Rata-rata persentase kesesuaian per kriteria dan keseluruhan dirangkum pada Tabel \ref{tab:ringkasan-llm}.

\begin{table}[H]
\centering
\caption{Ringkasan Rata-rata Persentase Kesesuaian Konten LLM}
\label{tab:ringkasan-llm}
\begin{tabular}{|L{7cm}|L{3cm}|}
\hline
\textbf{Kriteria} & \textbf{Persentase (\%)} \\
\hline
K1 - Kebenaran Konseptual & 100,00 \\
\hline
K2 - Kesesuaian Kurikulum & 100,00 \\
\hline
K3 - Kejelasan dan Tidak Menyesatkan & 91,11 \\
\hline
\textbf{Rata-rata Keseluruhan} & \textbf{97,33} \\
\hline
\end{tabular}
\end{table}

Ditinjau per topik, rata-rata persentase kesesuaian \textit{Graph} adalah 100\%, \textit{Linked List} 98\%, dan \textit{Tree} 94\%. Seluruh 15 soal referensi memperoleh skor K1 dan K2 penuh (3/3), yang berarti LLM tidak pernah memberikan penilaian yang secara konseptual salah maupun di luar cakupan kurikulum ASD pada seluruh kasus uji.

Penurunan skor hanya terjadi pada kriteria K3 (Kejelasan dan Tidak Menyesatkan) pada 4 dari 15 soal, yaitu LL-04, TR-01, TR-03, dan TR-04, dengan pola kesalahan yang serupa: LLM meminta elaborasi atau contoh tambahan yang sebenarnya berada di luar cakupan jawaban referensi (dan setara dengan tingkat kedalaman jawaban mahasiswa yang dinilai), sehingga \textit{feedback} yang diberikan terkesan lebih ketat/mendalam dibandingkan standar referensi, bukan karena mengandung kesalahan konsep atau miskonsepsi. Pola ini paling banyak ditemukan pada topik \textit{Tree} (3 dari 5 soal), sementara topik \textit{Graph} tidak menunjukkan kasus serupa sama sekali (skor sempurna pada seluruh soal).

Berdasarkan hasil tersebut, rata-rata persentase kesesuaian keseluruhan sebesar 97,33\% jauh melampaui indikator keberhasilan minimal 70\%, sehingga evaluasi kualitas konten LLM dinyatakan \textbf{lulus}. Konten yang dihasilkan oleh LLM, baik dalam bentuk soal latihan maupun evaluasi jawaban pengguna, terbukti relevan, akurat, dan konsisten dengan kurikulum ASD, dengan satu catatan bahwa tingkat "kekritisan" evaluasi LLM perlu dikalibrasi ulang agar tidak melebihi standar kedalaman yang diharapkan dari mahasiswa.

\section{Pembahasan Hasil Evaluasi}

Bagian ini membahas keterkaitan antar hasil dari ketiga tahap evaluasi yang telah dilakukan, yaitu pengujian \textit{black box}, pengujian kegunaan, dan evaluasi kualitas konten LLM, guna menjawab kedua rumusan masalah penelitian ini secara menyeluruh.

Hasil pengujian \textit{black box} menunjukkan bahwa seluruh 18 skenario uji dari 14 fitur fungsional dinyatakan lulus (Tabel \ref{tab:hasil-blackbox}), yang berarti seluruh kebutuhan fungsional yang ditetapkan pada Bab III, mulai dari alur \textit{guided learning}, mekanisme penguncian tab, integrasi LLM untuk \textit{generate} konten, hingga \textit{adaptive learning engine} berdasarkan \textit{weak concepts}, telah berjalan sebagaimana mestinya dari sisi fungsional. Namun demikian, pengujian \textit{black box} hanya memverifikasi kebenaran fungsional (apakah sistem menghasilkan keluaran sesuai spesifikasi) dan tidak menilai kualitas isi maupun pengalaman penggunaan keluaran tersebut, sehingga diperlukan dua evaluasi lanjutan yang saling melengkapi.

Pengujian kegunaan menunjukkan bahwa seluruh responden (5 dari 5) berhasil menyelesaikan seluruh tugas dengan tingkat keberhasilan 100\%, dan skor SUS rata-rata sebesar 93,0 berada pada kategori \textit{Excellent}, jauh di atas indikator keberhasilan minimal 68 (kategori \textit{Good}). Hal ini mengonfirmasi bahwa alur \textit{guided learning} yang telah diverifikasi lulus pada pengujian \textit{black box} juga dapat diikuti dengan mudah, efisien, dan memuaskan oleh mahasiswa dalam praktiknya. Meskipun demikian, wawancara pasca-tugas mengungkap sejumlah catatan perbaikan yang tidak tertangkap oleh pengujian \textit{black box} maupun tercermin pada skor SUS, seperti kebutuhan visualisasi tambahan pada materi struktur data, waktu tunggu \textit{generate} konten oleh LLM yang dirasa cukup lama, serta persepsi bahwa validasi jawaban oleh AI pada beberapa kasus terlalu ketat.

Evaluasi kualitas konten LLM menunjukkan rata-rata persentase kesesuaian keseluruhan sebesar 97,33\%, jauh melampaui \textit{threshold} 70\% yang ditetapkan sebagai indikator keberhasilan, dengan skor sempurna pada kriteria Kebenaran Konseptual (K1) dan Kesesuaian Kurikulum (K2). Temuan ini menjawab rumusan masalah kedua penelitian, yaitu bahwa LLM mampu menghasilkan konten pembelajaran (soal latihan dan evaluasi jawaban) yang relevan, akurat, dan konsisten dengan kurikulum ASD tanpa menimbulkan miskonsepsi yang signifikan. Penurunan skor yang teramati pada kriteria Kejelasan dan Tidak Menyesatkan (K3), khususnya pada topik \textit{Tree}, konsisten dengan pola yang sama, yaitu LLM cenderung mengevaluasi lebih ketat/mendalam dibandingkan standar jawaban referensi.

Keterkaitan yang menarik ditemukan antara hasil pengujian kegunaan dan evaluasi kualitas konten LLM: keduanya secara independen mengidentifikasi isu yang sama, yaitu kecenderungan LLM untuk melakukan validasi/evaluasi jawaban secara lebih ketat dari yang diharapkan pengguna maupun standar referensi (disebutkan langsung oleh salah satu responden pengujian kegunaan, dan teridentifikasi pada 4 dari 15 soal pada evaluasi kualitas konten). Konvergensi temuan dari dua metode evaluasi yang berbeda ini memperkuat validitas temuan tersebut dan menunjukkan bahwa kalibrasi tingkat "kekritisan" \textit{prompt} evaluasi LLM merupakan area perbaikan prioritas untuk pengembangan aplikasi selanjutnya, di luar aspek fungsional dan kebenaran konseptual yang sudah terbukti baik.

Secara keseluruhan, ketiga hasil evaluasi saling melengkapi dan menunjukkan bahwa aplikasi Teman Belajar ASD telah memenuhi seluruh kebutuhan fungsional (\textit{black box}), mudah dan nyaman digunakan oleh mahasiswa (\textit{usability}), serta menghasilkan konten pembelajaran dinamis yang akurat dan sesuai kurikulum (kualitas konten LLM). Evaluasi ini memiliki sejumlah keterbatasan yang perlu diperhatikan dalam menginterpretasikan hasil. Pertama, pengujian kegunaan hanya dilakukan pada satu topik (\textit{Linked List}) sehingga belum merepresentasikan pengalaman pengguna pada topik-topik lain dengan tingkat kompleksitas berbeda. Kedua, jumlah responden pengujian kegunaan (5 orang) memenuhi ambang minimal menurut Nielsen (1993) untuk mengidentifikasi sebagian besar masalah kegunaan, namun tetap tergolong kecil untuk generalisasi statistik yang kuat. Ketiga, evaluasi kualitas konten LLM dilakukan oleh satu evaluator (peneliti sendiri) terhadap 15 soal pada 3 dari 10 topik kurikulum, sehingga meskipun rubrik penilaian telah terstruktur, penilaian tetap mengandung unsur subjektivitas dan belum mencakup \textit{inter-rater reliability}. Cakupan topik yang dievaluasi juga belum merepresentasikan seluruh materi ASD yang diajarkan pada aplikasi.
